\documentclass[11pt]{article}
\usepackage[utf8]{inputenc}
\usepackage[T1]{fontenc}
\usepackage{amsmath,amsthm,amssymb,amsfonts}
\usepackage{geometry}
\usepackage{hyperref}
\geometry{margin=0.9in}
\setlength{\parskip}{2pt}

\newtheorem{theorem}{Theorem}
\newtheorem*{conjecture*}{Conjecture (Kislitsyn--Fredman--Linial)}

\newcommand{\BK}{\mathrm{BK}}
\newcommand{\LP}{\mathcal{L}(P)}
\newcommand{\Pp}{\mathbb{P}}

\title{\vspace{-2em}The $1/3$--$2/3$ Conjecture for Width-3 Posets:\\
       A Two-Page Summary}
\author{D.~Miller}
\date{\today}

\begin{document}
\maketitle
\vspace{-2em}

\paragraph{Statement.}
Let $P$ be a finite poset (not a chain) and let $\LP$ be the set of its
linear extensions under the uniform measure. A pair $(x,y)$ of
incomparable elements is \emph{balanced} if
$\Pp[x<_L y]\in[1/3,2/3]$ for $L\sim\mathrm{Unif}(\LP)$.

\begin{conjecture*}
Every non-chain finite poset has a balanced pair.
\end{conjecture*}

\noindent
Open since 1968. The current best unconditional constant is
$0.2764\ldots$ (Kahn--Linial 1991, refining Brightwell). It is known
for width-2 posets (Linial 1984), $2$-dimensional posets (Kahn--Saks
1984), semiorders (Brightwell), and symmetric posets, but
\emph{not} previously known at width~$3$.

\begin{theorem}[Main result]
Let $P$ be a finite poset of width at most~$3$ that is not a chain.
Then $P$ admits a pair $(x,y)$ of incomparable elements with
$1/3\le\Pp[x<_L y]\le 2/3$.
\end{theorem}

\section*{The two ideas}

The proof rests on two complementary observations. The first is a
\emph{global} tool that turns a question about one number (the
balance constant) into a question about the geometry of an
exponentially large graph. The second is a \emph{local} constraint
--- specific to width~$3$ --- that prevents that geometry from
accommodating an imbalanced poset.

\paragraph{Idea 1 — The balance constant is BK conductance.}
The Bubley--Karzanov graph $\BK(P)$ has vertex set $\LP$, with edges
joining extensions that differ by a single $P$-compatible adjacent
transposition. This is the natural local move on linear extensions;
Bubley--Karzanov (1998) showed the corresponding random walk mixes in
polynomial time.

The key reframing: suppose $(x,y)$ is the most-balanced incomparable
pair of an unbalanced $P$, with $\Pp[x<y]<1/3$. Consider the cut
$ S_{xy} = \{L\in\LP : x<_L y\}\subseteq\BK(P).$
Then (i) $S_{xy}$ has volume fraction $<1/3$, and (ii) every boundary
edge of $S_{xy}$ swaps $x$ and $y$ specifically, which is rare since
most extensions do not place $x,y$ adjacent. So a counterexample
produces a \emph{low-conductance cut} in $\BK(P)$. A Cheeger-type
inequality translates the original numerical statement
(``$\delta(P)<1/3$'') into a geometric one. Instead of arguing
about probabilities via correlation inequalities (Kahn--Saks,
Brightwell, Kahn--Linial), we now have an entire graph and an entire
cut to analyse, and the full strength of expansion methods becomes
available.

\paragraph{Idea 2 — Width 3 makes the obstructions essentially 2D.}
What could prevent a low-conductance cut from existing? The local
geometry of $\BK(P)$ near a generic ``rich'' incomparable pair
$(x,y)$ --- one sharing many common neighbors --- is a 2D grid:
fixing the relative order of all elements outside $\{x,y\}$ and their
joint neighborhood, the remaining freedom in $L$ is recorded by two
integer coordinates (``how far has $x$ advanced?'' / ``how far has
$y$ advanced?'') ranging over a region $D_{xy}\subseteq[0,t]^2$, and
BK moves within this fiber are unit grid moves.

This 2D coordinate system exists \emph{because} width~$\le 3$. A
fourth uninvolved element would push the local geometry to 3D, where
many more boundary configurations become possible; at width~$3$ the
only ``other'' element is one of the common neighbors, whose position
is recorded inside the grid.

The 2D picture has two consequences. \emph{(a) Local cuts are nearly
monotone.} A standard edge-isoperimetric stability theorem says: a
subset of a 2D grid region with small edge boundary is close to a
monotone staircase. The restriction of any low-conductance cut to a
rich fiber thus points in a definite direction --- assigning each
rich pair a canonical sign $\eta_{xy}\in\{+,-\}$ and axis. \emph{(b)
Two incoherent interfaces are incompatible.} If two rich pairs have
monotone interfaces pointing in incompatible directions, their joint
constraint on a shared fiber is over-determined: many BK boundary
edges are forced. This converts ``many rich pairs'' into a violation
of the conductance lower bound.

\paragraph{Putting it together.}
A width-$3$ counterexample $P$ would force --- by Idea~1 --- a
low-conductance cut $S$. By Idea~2(a), $S$ is locally monotone on
each rich fiber, giving each rich pair a canonical sign. By
Idea~2(b), if the signs are \emph{not} globally consistent, two
incoherent rich pairs alone exceed the boundary budget. So the signs
must be globally consistent --- which forces $P$ to admit a single
linear potential function aligned with $S$, collapsing $P$ to a
layered, essentially $1$-dimensional structure. But such posets are
already balanced (Linial's width-2 result and a small-$n$ enumeration
cover the residual cases). Contradiction.

\paragraph{Why width 3 is essential.}
Both ingredients use width~$3$. The 2D grid coordinatisation requires
that fixing two elements leaves no further incomparable third axis;
at width~$\ge 4$ the local geometry is at least 3D and grid
edge-isoperimetric stability is much weaker. The Dilworth dichotomy
producing ``enough rich interfaces'' also relies on a decomposition
into at most $3$ chains. Beyond width~$3$, both fronts need
genuinely new ingredients.

\paragraph{Repository.}
The full proof is in \texttt{main.tex} (assembling
\texttt{step1.tex}--\texttt{step8.tex}); a partial Lean~4
formalisation is in \texttt{lean/}; a self-audit and reviewer
pointers are in \texttt{summary.md} and \texttt{reviewers.md}. The
README is candid about authorship, AI assistance, and the absence of
external review --- please read it before citing.

\end{document}
